\subsection{Deriving the importance ratio}
\label{subsec:grpo-importance-ratio}

This subsection returns to the actual reward term $\rho_{i,t}A_i$ inside
Equation~\ref{eq:grpo-full-loss}. The probability ratio is not a second reward signal. It
follows from estimating an expectation under a new policy using actions sampled from an
old policy. Fix a
prompt-prefix state $s$ and treat the already computed advantage
$A_{\mathrm{old}}(s,a)$ as constant. The new policy's expected advantage is
\begin{align}
  J_s(\theta)
    &= \sum_a \pi_\theta(a\mid s)A_{\mathrm{old}}(s,a) \notag\\
    &= \sum_a \pi_{\mathrm{old}}(a\mid s)
       \frac{\pi_\theta(a\mid s)}{\pi_{\mathrm{old}}(a\mid s)}
       A_{\mathrm{old}}(s,a) \notag\\
    &= \mathbb{E}_{a\sim\pi_{\mathrm{old}}}
       \left[\rho_\theta(s,a)A_{\mathrm{old}}(s,a)\right],
  \qquad
  \rho_\theta(s,a)=\frac{\pi_\theta(a\mid s)}{\pi_{\mathrm{old}}(a\mid s)}.
  \label{eq:grpo-importance-weighted-expectation}
\end{align}
This is importance weighting, or a change of sampling distribution. The advantage supplies
the direction and magnitude of preference. The ratio only states how much probability
mass the new policy assigns to an old-policy sample. It is required wherever
$\pi_\theta$ is evaluated using data drawn from $\pi_{\mathrm{old}}$.

Consider two actions. Under the old policy, the good action has probability $0.25$ and
advantage $+3$; the bad action has probability $0.75$ and advantage $-1$. The old expected
advantage is zero. If the new policy assigns both actions probability $0.5$, its expected
advantage is one. Importance weighting the old distribution gives exactly the same result:
\begin{equation}
  0.25\left(\frac{0.5}{0.25}\right)(3)
  +0.75\left(\frac{0.5}{0.75}\right)(-1)=1.
  \label{eq:grpo-importance-weight-example}
\end{equation}
Without the ratios, the old samples would still average to zero and would not represent
the changed policy.

Advantage alone is detached, so $\nabla_\theta A=0$. In the actual GRPO reward term, the
trainable dependence enters through $\rho_\theta$ in
Equation~\ref{eq:grpo-importance-weighted-expectation}. When
$\pi_\theta=\pi_{\mathrm{old}}$, the ratio equals one but its gradient does not vanish:
\begin{equation}
  \nabla_\theta\rho_\theta
  =\rho_\theta\nabla_\theta\log\pi_\theta.
  \label{eq:grpo-ratio-gradient}
\end{equation}
Thus $\nabla_\theta(\rho A)=A\nabla_\theta\log\pi_\theta$ at the behavior policy. After
averaging over a trajectory, this is exactly the gradient in
Equation~\ref{eq:grpo-advantage-weighted-gradient}. Therefore the log-probability
expression in Section~\ref{subsec:grpo-gradient-direction} was a teaching surrogate for
the initial gradient, whereas
$\rho A$ is the quantity that actually appears in the GRPO loss. The two are not added
together, and GRPO does not multiply $\rho A$ by a separate cross-entropy loss.

Before clipping, the reward contribution for trajectory $i$ is therefore
\begin{equation}
  J_i^{\mathrm{reward}}(\theta)
  =\frac{1}{T_i}\sum_{t=1}^{T_i}\rho_{i,t}(\theta)A_i.
  \label{eq:grpo-unclipped-reward-objective}
\end{equation}
Section~\ref{subsec:grpo-clipped-objective} replaces each $\rho_{i,t}A_i$ with its clipped
version and restores the
reference-policy KL term, producing the two pieces inside the brackets of
Equation~\ref{eq:grpo-full-loss}.

For a complete response, exact change of distribution would use a trajectory ratio:
\begin{equation}
  \frac{\pi_\theta(o\mid q)}{\pi_{\mathrm{old}}(o\mid q)}
  =\prod_{t=1}^{T}
    \frac{\pi_\theta(o_t\mid q,o_{<t})}
         {\pi_{\mathrm{old}}(o_t\mid q,o_{<t})}.
  \label{eq:grpo-complete-trajectory-ratio}
\end{equation}
The product becomes high-variance for long sequences. GRPO instead holds the behavior
policy's sampled prefixes fixed and applies a local per-token ratio, broadcasting the
trajectory advantage across those token terms. This is a local policy-improvement
surrogate around $\pi_{\mathrm{old}}$, not exact reweighting after an arbitrarily large
policy change. The resulting per-token quantity is the $\rho_{i,t}$ used in the complete
loss in Equation~\ref{eq:grpo-full-loss}.
